% paper/sections/validation.tex — Validation and Correctness section
% Revised in the pre-arXiv review pass (cross-implementation agreement added).
\section{Validation and Correctness}
\label{sec:validation}

\subsection*{Development methodology disclosure}

\texttt{fdars} and the underlying \texttt{fdars-core} Rust crate were developed
with substantial assistance from Anthropic's Claude Code, an AI coding assistant.
All algorithms, architecture decisions, and final code were specified, reviewed,
and verified by the author; Claude Code served as an implementation accelerator
throughout the development lifecycle.  Because much of the implementation was
produced with AI assistance, independent validation is treated as a first-class
design constraint rather than an afterthought: every layer of the stack is
guarded by automated checks that run on each commit, results are compared with
independent implementations (Section~\ref{subsec:crossimpl}), and this
manuscript is itself a reproducibility artefact whose code snippets and figures
are regenerated and verified in CI.

\subsection{Agreement with independent implementations}
\label{subsec:crossimpl}

Unit tests show that code does what its author intended; they cannot show that
the author's intent matches the literature.  We therefore compare six
\texttt{fdars} methods with established, independently written R
implementations (Table~\ref{tab:crossimpl}).  The reference outputs were
computed once with \texttt{fda.usc} \citep{febrerobande_fdausc_2012},
\texttt{roahd} \citep{r_roahd}, \texttt{fdasrvf} \citep{r_fdasrvf}, and
\texttt{fdapace} \citep{r_fdapace} and committed together with their inputs;
the \texttt{fdars} side is recomputed on every CI run and a drift gate fails
if the table changes.

% Auto-generated by paper/code/crossimpl.py -- do not edit.
% References: paper/code/crossimpl_ref/ (produced by crossimpl_ref.R).
\begin{table}[htbp]
\centering
\footnotesize
\caption{Numerical agreement between \texttt{fdars} and independent R implementations. fdars is recomputed on every CI run against frozen reference outputs; $\rho$ is Spearman rank correlation.}
\label{tab:crossimpl}
\begin{tabular}{p{0.19\linewidth}p{0.27\linewidth}p{0.24\linewidth}p{0.22\linewidth}}
\toprule
\textbf{Method (data)} & \textbf{fdars vs reference} & \textbf{Agreement} & \textbf{Remark} \\
\midrule
Fraiman--Muniz depth (growth, $n=93$) & \texttt{depth.fraiman\_muniz\_1d} vs fda.usc~2.2.0 \texttt{depth.FM} & max $|\Delta|$ $<10^{-12}$; $\rho=1.000$ & exact; fdars \texttt{scale=True} $\equiv$ fda.usc \texttt{scale=FALSE} \\
\addlinespace
Modified band depth (growth) & \texttt{depth.modified\_band\_1d} vs roahd~1.4.3 \texttt{MBD} & $\rho=0.999$; max $|\Delta|=0.012$ & same ranking; small value differences \\
\addlinespace
Dense FPCA, 4 components (growth) & \texttt{Fdata.to\_pc} vs fda.usc~2.2.0 \texttt{fdata2pc} & eigenfunction $|\cos|$ 0.972--1.000; PC1 share 81.0\% vs 80.4\% & quadrature weighting differs on the non-uniform grid \\
\addlinespace
$L^2$ distance matrix (growth) & \texttt{metric.lp\_self\_1d} vs fda.usc~2.2.0 \texttt{metric.lp} & max rel.\ diff.\ 7.0\% & reference $=$ trapezoid rule (reproduced to $<10^{-12}$); fdars' rule is exact for quadratics on this grid \\
\addlinespace
Elastic Karcher mean (20 phase-shifted bumps) & \texttt{alignment.karcher\_mean} vs fdasrvf~2.5.0 \texttt{time\_warping} & rel.\ $L^2$ diff.\ 1.7\%; peak 0.949 vs 0.950 & cross-sectional mean peak is 0.511 \\
\addlinespace
PACE sparse FPCA (120 curves, 4--8 obs.\ each) & \texttt{pace\_fpca.pace\_fpca} vs fdapace~0.6.0 \texttt{FPCA} & eigenfunction $|\cos|$ 0.993, 0.992; $\hat\lambda=(3.96, 0.86)$ vs $(3.69, 0.99)$ & true $\lambda=(4, 1)$; same Gaussian kernel and bandwidth, but fdars smooths local-constant, fdapace local-linear \\
\bottomrule
\end{tabular}
\end{table}


Fraiman--Muniz depth agrees to machine precision.  Modified band depth
produces the same ranking with value differences of at most $0.012$.  The
elastic Karcher mean matches \texttt{fdasrvf} to within 1.7\% in relative
$L^2$ norm, and both recover the sharp common peak that the cross-sectional
mean flattens.  The remaining differences trace to documented modelling
choices rather than errors.  The $L^2$ distances differ by up to 7\% because
\texttt{fda.usc} integrates with the trapezoid rule, which we reproduce
exactly, whereas \texttt{fdars} uses a higher-order composite rule that
integrates quadratics exactly on the non-uniform growth grid.  The same
quadrature weighting explains the small FPCA differences.  For PACE, both
libraries use a Gaussian kernel with the same bandwidth, but \texttt{fdars}
smooths the mean with a local-constant (Nadaraya--Watson) estimator where
\texttt{fdapace} uses local-linear smoothing; the eigenfunctions nevertheless
agree closely ($|\cos| \geq 0.99$) and both eigenvalue estimates are close to
the true values.

\subsection{Rust layer}

The numerical algorithms reside in the \texttt{fdars-core} crate, which ships
3{,}260 Rust unit and integration tests.  The PyO3 binding crate
(\texttt{pyfda}) contains no Rust unit tests of its own---it is tested from
Python---and CI enforces its code quality: \texttt{cargo fmt --check} checks
formatting, \texttt{cargo clippy -- -D warnings} promotes warnings to errors,
and \texttt{cargo test} builds the crate on both the current stable toolchain
and the declared minimum supported Rust version (1.83).

\subsection{Python test suite}

The \texttt{tests/} directory (5{,}700 tests) is exercised by \texttt{pytest}
on every CI run across Python~3.9 through~3.14.
Tests cover the \texttt{Fdata} container, numerical operations (smoothing,
depth, alignment, regression, monitoring, covariance sampling), the advisor
and MCP server, and the dataset loaders.

\subsection{Property-based numerical tests}

\texttt{tests/test\_r\_parity.py} checks mathematically verifiable properties
of the outputs---integrals of known functions, determinism under fixed seeds,
correct output shapes, and expected statistical behaviour (for example, that a
CUSUM chart detects a known mean shift).  Despite the file name, these tests do
not compare against R; that comparison is the one reported in
Table~\ref{tab:crossimpl}.  They span depth measures, $L^2$ inner products,
Simpson integration, landmark registration, nonparametric regression,
functional monitoring (CUSUM and EWMA), elastic changepoint detection, and
prediction metrics.

\subsection{scikit-learn API compliance}

\nsklearnestimators{} estimators in \texttt{fdars.sklearn} pass the
scikit-learn \texttt{check\_estimator} conformance suite
\citep{pedregosa_sklearn_2011}, run from \texttt{tests/sklearn/} in CI across
all supported Python versions.
Passing \texttt{check\_estimator} means that the estimators honour
scikit-learn's fit--predict contract, are picklable, and compose correctly
inside \texttt{Pipeline} and \texttt{GridSearchCV}.

\subsection{Reproducibility as validation}

This manuscript is itself a validation artefact.
Every code snippet in Section~\ref{sec:capabilities}, every figure, and every
number in Section~\ref{sec:casestudies} and Table~\ref{tab:crossimpl} is
produced by a committed \texttt{paper/code/} script.
CI drift gates verify that regenerating the snippets, the coverage macros, the
bibliography, and the agreement table leaves no diff against the committed
sources, and that the figures regenerate byte-identically.
No printed example can silently diverge from the live API.
